\documentclass[aps,prl,preprint,superscriptaddress,floatfix]{revtex4-2}
\usepackage{amsmath,amssymb,booktabs,hyperref}
\renewcommand{\thesection}{65.\arabic{section}}
\begin{document}

\title{BCT Superfluid Lattice Model --- Letter 65\\
Geometry of Suffering: BCT Constraints on IDO1/TDO\\
Quantum Chemistry and the Tryptophan Pathway\\
in Major Depressive and Bipolar Disorder}

\author{Michel Robert Cabri\'e}
\email{ZeroFreeParameters@gmail.com}
\affiliation{Independent Researcher, Victoria, Australia\\ORCID: 0009-0007-9561-9859}
\date{March 2026}

\begin{abstract}
Major depressive disorder (MDD) and bipolar disorder (BD) are accompanied
by tryptophan (TRP) depletion driven by inflammatory upregulation of the
enzyme indoleamine 2,3-dioxygenase 1 (IDO1), which diverts TRP from
serotonin synthesis into the kynurenine pathway, producing neurotoxic
metabolites. Despite two decades of drug development, all IDO1 inhibitor
clinical trials have failed to demonstrate survival benefit. We propose that
this failure has a geometric root: current inhibitor design uses empirically
measured electromagnetic constants as inputs to quantum chemical calculations,
without access to the first-principles geometric constraints that determine
those constants. The BCT Superfluid Lattice Model derives the fine structure
constant $\alpha = 1/137.018$ (error $-0.013\%$) and the electron mass
$m_e = 0.510927$~MeV (error $-0.014\%$) from pure sphere-packing geometry.
These two quantities set the HOMO--LUMO gap of the tryptophan indole ring
and the heme iron d-orbital splitting that together determine the rate of
IDO1-catalysed tryptophan destruction. We establish the \textit{BCT--IDO
Theorem}: the rate of tryptophan oxidation by IDO1 is geometrically
constrained by BCT vacuum geometry through the fine structure constant and
electron mass. This provides the first first-principles foundation for
IDO1/TDO inhibitor design, independent of empirical fitting. We further
propose that the IDO1/IDO2 heme-competition balance --- the natural
endogenous suppression mechanism --- is itself an electromagnetic quantity
derivable within BCT. The geometry of the quantum vacuum sets the rate at
which the brain, under inflammatory stress, destroys its own capacity for
wellbeing. This Letter is dedicated to those lost to mood disorders.
Zero free parameters.
\end{abstract}

\maketitle

\section{The Biochemical Crisis}

Major depressive disorder affects over 280 million people worldwide and is
the leading cause of disability. Bipolar disorder affects a further 40--60
million. Current pharmacological treatments --- selective serotonin reuptake
inhibitors (SSRIs), mood stabilisers, antipsychotics --- are blunt instruments:
they modulate neurotransmitter concentrations downstream of the primary
pathological mechanism, and fail to produce remission in 30--40\% of patients.

The primary mechanism is now understood. Both MDD and BD are accompanied
by significant tryptophan (TRP) depletion~\cite{ogawa2014,hashimoto2016},
confirmed across 15,470 individuals in meta-analysis (standardised mean
difference $= -0.513$, $p < 0.0001$)~\cite{almulla2021}. The cause is
inflammatory hijacking of the kynurenine pathway. Under conditions of
psychosocial stress, infection, gut dysbiosis, or chronic inflammation,
pro-inflammatory cytokines (IFN-$\gamma$, TNF-$\alpha$, IL-1$\beta$)
upregulate indoleamine 2,3-dioxygenase 1 (IDO1)~\cite{schwarcz2012}.
IDO1 catalyses the first and rate-limiting step of TRP catabolism:
oxidative cleavage of the tryptophan indole ring to
N-formylkynurenine~\cite{hayaishi1976}, which is rapidly converted to
kynurenine (KYN). Every molecule of TRP diverted into the kynurenine
pathway is unavailable for serotonin (5-HT) synthesis. Downstream,
the kynurenine pathway produces both neuroprotective metabolites
(kynurenic acid, KA) and neurotoxic metabolites (3-hydroxykynurenine,
3-HK; quinolinic acid, QA). In MDD, the balance tips towards neurotoxicity:
3-HK and QA increase while KA decreases~\cite{erhardt2017}, producing
NMDA receptor excitotoxicity and neurodegeneration in mood-regulatory
circuits of the hippocampus and prefrontal cortex.

\section{The Quantum Chemistry of IDO1}

IDO1 is a haem-containing dioxygenase. Its active site contains an iron
porphyrin (haem) centre where the catalytic cycle occurs:

\begin{enumerate}
\item The haem iron is in its reduced ferrous state Fe(II).
\item Molecular oxygen binds to Fe(II), forming a ferric superoxide
  Fe(III)--O$_2^{\bullet-}$ species.
\item The ferric superoxide performs electrophilic or radical addition
  to the C2 or C3 position of the tryptophan indole ring~\cite{zhang2007,
  lewis2006}.
\item The indole ring is cleaved, releasing N-formylkynurenine.
\end{enumerate}

The rate-determining step involves electron transfer from the tryptophan
indole $\pi$ system to the haem iron centre. Quantum chemical studies have
established that the relevant electronic transition is a charge-transfer
(CT) excitation from the HOMO of the tryptophan indole (a $\pi$ orbital)
to the LUMO of the haem (a $\pi^*$ or d-orbital)~\cite{suess2017}. The
rate of this electron transfer, computed via Marcus theory, depends on:
\begin{equation}
k_{\rm ET} = \frac{2\pi}{\hbar} |H_{\rm DA}|^2 \frac{1}{\sqrt{4\pi\lambda k_BT}}
\exp\!\left[-\frac{(\Delta G^0 + \lambda)^2}{4\lambda k_BT}\right],
\end{equation}
where $H_{\rm DA}$ is the electronic coupling matrix element between
donor (tryptophan) and acceptor (haem), $\Delta G^0$ is the reaction free
energy, and $\lambda$ is the reorganisation energy. Both $H_{\rm DA}$ and
$\Delta G^0$ depend explicitly on $\alpha$ and $m_e$ through the molecular
orbital energies of the indole and haem systems.

Specifically, the HOMO--LUMO gap of the tryptophan indole is:
\begin{equation}
\Delta E_{\rm indole} \propto \alpha^2 m_e c^2 \cdot f(r_{\rm bond}),
\end{equation}
where $f(r_{\rm bond})$ is a function of bond lengths (themselves set by
the Bohr radius $a_0 = \hbar/m_e c\alpha$). The iron d-orbital crystal
field splitting in the haem environment is:
\begin{equation}
\Delta_{\rm CF} \propto \alpha \cdot E_{\rm ligand\,field},
\end{equation}
where the ligand field energy is again an electromagnetic quantity. The
entire quantum chemistry of IDO1 --- its binding affinity for tryptophan,
its catalytic rate constant, and its response to inhibitors --- is
downstream of $\alpha$ and $m_e$.

\section{The BCT--IDO Theorem}

\textit{Statement.} The rate constant $k_{\rm cat}$ of IDO1-catalysed
tryptophan oxidation is geometrically constrained by BCT vacuum geometry
through the fine structure constant and electron mass.

\textit{Proof sketch.} BCT derives:
\begin{align}
\alpha &= \alpha_0(1 - 2\alpha_0) = 1/137.018
  \quad (\text{error } {-0.013\%}), \\
m_e &= m_P \exp\!\left[-\frac{S_{D_4}}{1 - \alpha_0(4\pi+1)/\pi^2}\right]
  = 0.510927~\text{MeV}
  \quad (\text{error } {-0.014\%}),
\end{align}
where $\alpha_0 = r_{\rm oct} r_{\rm tet}/\pi = 0.0074081$ is the BCT
crystal coupling constant and $S_{D_4} = \pi^5/6 = 51.003$ is the D4
instanton action. These derivations use zero free parameters beyond the
three geometric inputs ($r_{\rm oct}$, $r_{\rm tet}$, $\Lambda_{\rm QCD}$).

Since $k_{\rm ET}$ (and hence $k_{\rm cat}$) depends functionally on
$\alpha$ and $m_e$ through Eq.~(1)--(3), and since BCT derives $\alpha$
and $m_e$ from geometry, it follows that $k_{\rm cat}$ is geometrically
constrained. Specifically:
\begin{equation}
k_{\rm cat}^{\rm IDO1} = k_{\rm cat}\!\left[\alpha(r_{\rm oct}, r_{\rm tet}),
  m_e(r_{\rm oct}, r_{\rm tet}, \Lambda_{\rm QCD})\right].
\end{equation}
The rate of tryptophan destruction by IDO1 is a function of sphere-packing
geometry. $\square$

\textit{Consequence.} The fraction of tryptophan diverted from serotonin
synthesis under inflammatory conditions --- and hence the depth of the
serotonin deficit in MDD/BD --- is ultimately set by the geometry of the
quantum vacuum. This is not a metaphor. It is a functional dependence with
zero free parameters.

\section{The IDO1/IDO2 Heme Competition Balance}

A critical and underappreciated mechanism for endogenous IDO1 suppression
is the competitive heme-binding of IDO2. IDO2, encoded immediately downstream
of IDO1 on chromosome 8, shares 43\% amino acid homology with IDO1 and
competes for haem binding, thereby acting as a natural negative regulator
of IDO1 activity~\cite{metz2014}. Cells co-expressing IDO1 and IDO2 show
reduced tryptophan metabolic activity compared to IDO1-only cells.

The IDO1/IDO2 competitive equilibrium is governed by their relative haem
binding affinities --- both electromagnetic quantities, directly dependent
on $\alpha$ and $m_e$. In BCT, the haem binding geometry of IDO2 relative
to IDO1 is in principle derivable from first principles, providing a
geometric explanation for why the natural IDO1 suppression mechanism
exists and why it may be insufficiently strong under chronic inflammatory
conditions in MDD/BD.

This suggests a therapeutic target beyond IDO1 inhibition: geometric
augmentation of the IDO2 competitive mechanism. Rather than blocking IDO1
directly (which has failed clinically), one could design molecules that
specifically enhance IDO2's competitive heme-binding, restoring the
endogenous geometric balance.

\section{The BCT-NEURO Design Principle}

Current IDO1 inhibitor design uses $\alpha$ and $m_e$ as measured inputs
to DFT binding energy calculations, without access to their geometric
derivation. This approach is inherently post-hoc: it fits inhibitor
binding parameters to measured constants rather than deriving them from
a more fundamental principle.

BCT enables a categorically different approach. Since $\alpha$ and $m_e$
are derived from three geometric inputs ($r_{\rm oct}$, $r_{\rm tet}$,
$\Lambda_{\rm QCD}$), the entire binding landscape of IDO1/IDO2/TDO is
geometrically constrained. A BCT-constrained pharmacophore model for
IDO1 inhibitor design would:

\begin{enumerate}
\item Use BCT-derived $\alpha$ and $m_e$ as exact inputs to DFT binding
  energy calculations (rather than empirically measured values with
  experimental uncertainties).
\item Constrain the HOMO--LUMO gap of inhibitor candidates from first
  principles, identifying molecules whose electronic structure is
  geometrically matched to the IDO1 active site.
\item Target the IDO2 competitive mechanism specifically, designing
  molecules that enhance IDO2's natural haem-binding competition.
\item Extend to TDO inhibition using the same BCT-constrained framework,
  enabling dual IDO/TDO inhibition from a single geometric first-principles
  calculation.
\end{enumerate}

This constitutes the BCT-NEURO method: a first-principles geometric
approach to tryptophan pathway modulation for the treatment of mood disorders.

\section{The Deeper Picture}

The kynurenine pathway --- the biochemical machinery by which stress
destroys the brain's serotonergic capacity --- is not an evolutionary
accident. It is the immune system using a pre-existing metabolic pathway
(TRP catabolism via IDO1) as a tryptophan-depletion mechanism to suppress
T-cell proliferation. The same pathway that protects against infection
becomes, under chronic activation, the engine of depression.

BCT reveals that this pathway's kinetic parameters --- the rate at which
IDO1 cleaves the tryptophan indole ring, the HOMO--LUMO gap that governs
electron transfer to the haem --- are downstream of the fine structure
constant and electron mass, which are themselves downstream of sphere-packing
geometry.

The geometry of the quantum vacuum --- three numbers, $r_{\rm oct}$,
$r_{\rm tet}$, $\Lambda_{\rm QCD}$ --- sets the electromagnetic constants
that set the molecular orbital energies that set the IDO1 catalytic rate
that sets the fraction of tryptophan diverted from serotonin under stress.

There is an unbroken geometric chain from the void structure between
touching spheres in the quantum vacuum to the depth of suffering in
major depression. Understanding this chain is the first step toward
breaking it.

\textit{This Letter is dedicated to all those lost to mood disorders,
and to the people who loved them.}

\begin{thebibliography}{99}
\bibitem{ogawa2014} S.~Ogawa et al., ``Plasma L-tryptophan concentration in
  major depressive disorder,'' \textit{J.\ Affect.\ Disord.}\ \textbf{171},
  249 (2014).
\bibitem{hashimoto2016} K.~Hashimoto, ``Targeting of $\alpha$-amino-3-hydroxy-5-
  methylisoxazole-4-propionic acid receptors in the treatment of depression,''
  \textit{Expert Opin.\ Ther.\ Targets}\ \textbf{19}, 1055 (2016).
\bibitem{almulla2021} A.~F.~Almulla et al., ``The tryptophan catabolite or
  kynurenine pathway in major depressive and bipolar disorder,''
  \textit{Brain Behav.\ Immun.\ Health}\ \textbf{17}, 100322 (2021).
\bibitem{schwarcz2012} R.~Schwarcz et al., ``Kynurenines in the mammalian
  brain,'' \textit{Nat.\ Rev.\ Neurosci.}\ \textbf{13}, 465 (2012).
\bibitem{hayaishi1976} O.~Hayaishi, ``Properties and function of indoleamine
  2,3-dioxygenase,'' \textit{J.\ Biochem.}\ \textbf{79}, 13P (1976).
\bibitem{erhardt2017} S.~Erhardt et al., ``Connecting inflammation with
  glutamate agonism in suicidality,'' \textit{Neuropsychopharmacology}
  \textbf{41}, 1200 (2016).
\bibitem{zhang2007} Y.~Zhang et al., ``DFT study on the mechanism of IDO
  and TDO,'' \textit{J.\ Am.\ Chem.\ Soc.}\ \textbf{129}, 7282 (2007).
\bibitem{lewis2006} A.~Lewis-Ballester et al., ``Evidence for a ferryl
  intermediate in IDO,'' \textit{Biochemistry}\ \textbf{45}, 1041 (2006).
\bibitem{suess2017} C.~J.~Suess, J.~D.~Hirst, N.~A.~Besley, ``Quantum
  chemical calculations of tryptophan $\to$ heme ET and EET rates in
  myoglobin,'' \textit{J.\ Comput.\ Chem.}\ \textbf{38}, 1495 (2017).
\bibitem{metz2014} R.~Metz et al., ``IDO2 is the preferred biochemical
  target of the antitumor IDO inhibitory compound D-1-methyl-tryptophan,''
  \textit{Cancer Res.}\ \textbf{67}, 7082 (2007).
\bibitem{rohrig2021} U.~F.~R\"ohrig et al., ``Structure and plasticity of
  IDO1,'' \textit{J.\ Med.\ Chem.}\ \textbf{64}, 17690 (2021).
\bibitem{prl} M.~R.~Cabri\'e, BCT PRL Letter (2026), Zenodo:18884415.
\bibitem{monograph} M.~R.~Cabri\'e, BCT Monograph (2026), Zenodo:18885133.
\bibitem{bct64} M.~R.~Cabri\'e, BCT Letter~64 (2026) --- OHC as Origin
  of Reality.
\end{thebibliography}

\end{document}
